\scanpage{31}
\uline{Lemme} $\unsure{\mathscr T|\mathscr E}$ un topos \inserted{loc.} annelé

$C$ Catégorie des couples \unsure{$E\sqrt{\vphantom G G}$|$E,G$} de modules loc. libres de t.f. \inserted{partout $\ne0$}

\margin{(tous les isos)}

$C'$ \rule{3em}{.35pt} quadruples $A,E,F,\varphi$,

\[
\left\{\begin{aligned}
&A\ \text{Algèbre sur }\mathscr T\ \text{loc. isom. à une }\inserted{\text{Alg. de matrices}}\\
&E,F\ \text{Mod. loc. lib. de t.f.}\ \inserted{A,E,F\ \text{partout }\ne0}\\
&\varphi\ \text{un isom. } A\otimes\underline{\mathrm{End}}\,E\simeq\underline{\mathrm{End}}(F).\\
&\hspace{1em}(\text{tous les isos})
\end{aligned}\right.
\]

$F:C\longrightarrow C'$ le foncteur

\[
F:(G,E)\longmapsto\underline{\mathrm{End}}(G),E,G\otimes E,\varphi_{G,E}
\]

où $\varphi_{E,G}$ est l'isom. \uline{can.} $\underline{\mathrm{End}}(G)\otimes\underline{\mathrm{End}}(E)\simeq\underline{\mathrm{End}}(G\otimes E)$

Alors $F$ est une équivalence

\uline{Dém} $F$ provient d'un hom de champs $\underline C\mathbin{\unsure{-|\to}}\underline C'$, et il suffit
de prouver que celui-ci est une équivalence. Or, si $\unsure{g,e|g,f|\mathcal G,\mathcal E}$ sont
deux entiers $\unsure{>|\geq}0$ \inserted{(resp. $a,e,f$ trois entiers $\geq0$)}, soit $\underline C_{g,e}$ le sous-champ de $\underline C$ formé des $E,G$ tels
que rang $E=\unsure{e|g}$, rang $G=g$

(resp. $\underline C'_{a,e,f}$ le sous-champ de $\underline C$ formé des quadruples
pour lesquels $A$ de rg $a^2$, $E$ de rg $e$, $F$ de rg $\unsure{af|f|ae}$)

On a

\[
\underline C_{\deleted{\unsure{g}}}\simeq\coprod_{g,e}\underline C_{\unsure{g}e}\qquad
\underline C'\simeq\coprod_{a,e,f}\underline C'_{a^{\unsure{2}},e,f}
\]

(\uline{somme de champs}) et $\underline F$ induit

\[
\deleted{\unsure{F_{g,e}}}\underline F_{g,e}:\underline C_{g,e}\longrightarrow\underline C'_{g^{2},e,}
\]

on est donc ramené à prouver que les $\underline F_{g,e}$ sont des équivalences.

\uline{Notons} $\underline C_{g,e}$ est une \uline{gerbe} \inserted{trivialisée \deleted{\unsure{liée}} par $\deleted{\unsure{M(\mathcal O_{\mathscr T})}},\ \xi=(\underline O_{\mathscr T}^{g},\underline O_{\mathscr T}^{\unsure{e}}\deleted{\unsure{G}})$} liée par $\underline{\mathrm{GL}}(e)\times\underline{\mathrm{GL}}(g)\simeq\underline{\mathrm{Aut}}(\xi)$

\deleted{$\unsure{\underline C_{g,e}}$} $\underline C'_{g,e}$ est une \uline{gerbe} \deleted{(\unsure{liée par} $\underline{\mathrm{GL}}(e)\times\underline{\mathrm{GL}}(g)$ ?)}
\inserted{(trivialisée par $\xi'=\varphi(\xi)$, \deleted{$\unsure{\mathrm{End}(O_{\mathscr T}^{g}),O_{\mathscr T}^{e},O_{\mathscr T}^{g}\otimes O_{\mathscr T}^{e},\varphi}$}}

\deleted{En effet, $\unsure{\underline C'_{g,e}|\underline C_{g,e}}$ loc (\unsure{et m\textsuperscript{me} globalmt}) a une section}

En effet, \unsure{l'est|c'est} un champ en groupoïdes, \deleted{donc} reste à prouver que
deux objets sont loc. isom. On peut supposer $A=A'$, $E=E'$, $F=F'$,
On considère deux isom $\varphi,\varphi':A\otimes\underline{\mathrm{End}}(E)\rightrightarrows\underline{\mathrm{End}}(F)$
Ils diffèrent par un autom. \deleted{\unsure{intérieur}} de $\underline{\mathrm{End}}(F)$, \deleted{d'où} or tt autom
de $F$ est (localement) intérieur, donc on gagne. \deleted{De plus,}

Pour terminer la démonstration, il reste à prouver
que $\underline F_{g,e}$ induit $\underline{\mathrm{Aut}}(\xi)\xrightarrow{\alpha}\underline{\mathrm{Aut}}(\varphi(\xi))$ qui est un isom.

Or $\xi=(G,E)$ $\xi'=\underline{\mathrm{End}}(G),E,G\otimes E,\varphi_{G,E}:\underline{\mathrm{End}}(G)\otimes\underline{\mathrm{End}}(E)\simeq\underline{\mathrm{End}}(G\otimes E)$

Or si $u,v$ est un autom de $\xi$, $F(u,v)$ est l'aut
$\mathrm{End}(u),v,u\otimes v$, donc si $F(u,v)$ est connu, $v$ est connu et
$u\otimes v$ est connu, d'où \unsure{facilmt|facilemt} (comme $E$ partout non nul) $u$ connu,
donc $\alpha$ injectif. Soit maintenant $\sigma=(\lambda,v,w)$ un autom
% archive sheet number: 30
